\scanpage{13}
6. Pic-catégories.

On peut définir les Pic-catégories comme éléments d'une 3-catégorie, qui soit une sous-3-catégorie (pleine \inserted{\deleted{\unsure{à}} au niveau des 2-objet et 3-objet)} de celle des Gr-catégories. \deleted{\unsure{Ainsi|Et si|Soient} $\unsure{si}\ \unsure{\mathscr P|P},\mathscr P'$ sont deux Pic-catégories, et $\Lambda,\Lambda'\colon\mathscr P\rightrightarrows\mathscr P'$ deux $\otimes$ foncteurs \unsure{AC|AU}} Mais on peut aussi \unsure{penser|passer} à la 2-catégorie associée en gardant les $=\mathrm s$ 0 et 1 objets, mais en prenant les 2-objets à isom. près \inserted{(c'est en fait la relation triviale)}; \deleted{\unsure{on}} et \inserted{oubliant les 3-objets} se restreint aux cat. de Picard strictes, on retrouve alors la 2-catégorie de \uline{Deligne} (où les catégories \uline{Hom} sont elles-m$\overline{\rm m}$ des catégories de Picard déduites des $R\operatorname{Hom}(L^{\displaystyle\cdot},L'^{\displaystyle\cdot})$ par troncature).

Si on n'insiste pas sur les Pic-cat. \uline{strictes}, on peut dire ceci :

a) Soit $\mathscr P$ une Gr-catégorie. Elle admet une structure de Pic-catégorie compatible avec \unsure{le reste|la vente|la préd\textsuperscript{te}} (i.e. une \deleted{structure} contrainte de comm. compatible avec \unsure{le reste|la contr.}) ssi elle est splittable, i.e. $k(\mathscr P)\in H^3(\overset{\substack{M\\ \Vert}}{\pi_0},\overset{\substack{N\\ \Vert}}{\pi_1})$ est nulle Alors $\mathscr P$ admet $=$ une structure de Picard stricte.

b) Les structures de Picard \deleted{strictes} sur $\mathscr P$ forment un torseur sous $\operatorname{Bil}\deleted{\unsure{\mathrm{alt}|\mathrm{ant}}}\unsure{\mathrm{A}|\mathrm{Alt}|\mathrm{ant}|\mathrm{alt}}(\pi_0,\pi_0;\pi_1)$. Les structures de Picard strictes forment un sous-torseur sous $\operatorname{Bil\,alt}(\pi_0,\pi_0;\pi_1)$.
% archive sheet number: 12
